\section{Connect validation to publication}
The complete entry \texttt{png\_decode} takes the PNG address in \texttt{\$a0}, file
count in \texttt{\$a1}, framebuffer address in \texttt{\$a2}, and a context address
in \texttt{\$a3}. Its vector counterpart is \texttt{png\_decode\_msa}. The context
holds six caller inputs and three results, all native doublewords. It is a description
of caller-provided storage, not an instruction to allocate that storage.

\Needspace{14\baselineskip}
\begin{longtable}{p{22mm}p{44mm}p{84mm}}
  \toprule Byte offset & Field                 & Meaning                                              \\\midrule
  0                    & Destination stride    & Words between framebuffer row starts                 \\
  8                    & Framebuffer capacity  & Writable bytes beginning at the framebuffer argument \\
  16                   & Compressed workspace  & Address for joined IDAT payloads                     \\
  24                   & Compressed capacity   & Writable bytes at that address                       \\
  32                   & Row workspace         & Address for inflated and reconstructed rows          \\
  40                   & Row capacity          & Writable bytes at that address                       \\
  48                   & Width result          & Pixels per row, published on success                 \\
  56                   & Height result         & Row count, published on success                      \\
  64                   & Published byte result & Active pixel bytes, $4wh$                            \\\bottomrule
\end{longtable}
Keep the file, both workspaces, framebuffer, context and private stack storage
nonoverlapping. Unaligned byte buffers are permitted; the context requires eight-byte
alignment and the framebuffer requires four-byte alignment. Addresses retain their
high bits even when every selected count is small. The ordinary display stride is
512 words, while a tightly packed test buffer may use the image width.

Collect into compressed workspace and retain the returned dimensions privately.
Compare the predicted inflated count with row capacity before calling the inflater.
After exact decompression, reconstruct rows in that same workspace. Then pack their
samples into the framebuffer. Each failed stage returns to the caller without
publishing context results. Scratch contents may have changed, but a returned format
or capacity failure leaves the framebuffer and all context bytes unchanged.

On success, write width, height and active byte count to the context, and return zero
in \texttt{\$v0} with the active count in \texttt{\$v1}. On failure, return one and
zero. Retain the input file, context inputs, row padding and saved machine state.
This order makes the final-filter test useful: even with valid outer CRCs and a valid
Adler checksum, an unsupported filter selector must stop before any pixel publication.
